%% LyX 2.0.7.1 created this file.  For more info, see http://www.lyx.org/.
%% Do not edit unless you really know what you are doing.
\documentclass[11pt,english]{article}
\usepackage[latin9]{inputenc}
\usepackage{geometry}
\geometry{verbose,tmargin=1in,bmargin=1in,lmargin=1in,rmargin=1in}
\usepackage{varioref}
\usepackage{float}
\usepackage{textcomp}
\usepackage{amstext}
\usepackage{graphicx}
\usepackage{esint}

\makeatletter

%%%%%%%%%%%%%%%%%%%%%%%%%%%%%% LyX specific LaTeX commands.
\DeclareRobustCommand{\greektext}{%
  \fontencoding{LGR}\selectfont\def\encodingdefault{LGR}}
\DeclareRobustCommand{\textgreek}[1]{\leavevmode{\greektext #1}}
\DeclareFontEncoding{LGR}{}{}
\DeclareTextSymbol{\~}{LGR}{126}
\newcommand{\lyxmathsym}[1]{\ifmmode\begingroup\def\b@ld{bold}
  \text{\ifx\math@version\b@ld\bfseries\fi#1}\endgroup\else#1\fi}

%% Because html converters don't know tabularnewline
\providecommand{\tabularnewline}{\\}
%% A simple dot to overcome graphicx limitations
\newcommand{\lyxdot}{.}


\makeatother

\usepackage{babel}
\begin{document}

\title{Basic Interferometer on Point Sources, the Sun, and the Moon}


\author{By Patrick Kantorski}


\date{April 8, 2014}
\maketitle
\begin{abstract}
This report was written for a radio astronomy undergraduate course
at UC Berkeley. Radio interferometry is a powerful observational method
used in Radio Astronomy for observing the sky in the Fourier domain.
In this lab, my partners and I used a local radio interferometer at
Berkeley to observe the Sun, the Moon, and point sources. We were
also able to approximate the base line of our interferometer and measure
the diameter of the Sun and the Moon.
\end{abstract}

\section{Introduction}

The purpose of this lab on interferometry was to demonstrate diffraction
theory as it applies to a real radio interferometer. Fringes along
with their amplitudes and phases were examined along with the ``Fourier
transform'' relationship between the interferometer's response and
luminosity of the sky. Additionally, my partners and I had to learn
a lot about time, coordinates, and conversions to progress in this
lab. We also obtained horizon-to-horizon data on various astronomical
sources ranging from point sources to the Sun and Moon. When performing
data analysis, we use linear and nonlinear least-squares fitting methods
to filter out accurate point source positions and obtain approximate
diameters for the Sun and Moon along with an estimate value of our
two-element interferometer's baseline.


\section{Background Theory}


\subsection{Astronomical Coordinates}

There are many different ways to go about keeping track of time and
position using coordinate systems. For us to make and record accurate
astronomical observations, we must know both where on Earth we are
positioned and the current time relative to such location. By knowing
both the position and time of observations, the direction from which
a telescope is pointed can be accurately deciphered. Latitude and
longitude coordinates work well to encode relative positions on the
Earth. However because the Earth is not perfectly spherical, detailed
surveys using latitude and longitude are not quite as simple as they
may appear. On the contrary, accurately keeping track of time has
its own issues. On Earth, clocks are tuned to different times according
to which time zone and year a person is currently in. However in astronomy,
using time zones to keep track of time is not particularly useful.
To make counting time easier, astronomers use systems that do away
with time zones and synchronize time to the ``solar day'' at a position
of 0\textdegree{} longitude, which also happens to conveniently run
through the British Royal Observatory.

\begin{figure}[H]
\begin{centering}
\includegraphics[scale=0.5]{/Users/ppkantorski/Documents/Radio_Astronomy/Lab_3/Pictures/sidereal_day}
\par\end{centering}

\caption{On Earth, because the planet experiences prograde motion, a sidereal
day is slightly shorter than a solar day. At position 1, the Sun and
a distant star behind are both directly overhead. At position 2, the
Earth rotates about 360\textdegree{}, making the distant star behind
the sun appear again directly overhead. This shift from position 1
to position 2 is called a sidereal day. However, by the time the Earth
rotates back to face the Sun, as seen in position 3, the Sun is again
directly overhead, but the distant star behind is not. Since the revolution
around the Sun changes the angle from which the Sun is seen relative
to the Earth, position 1 to position 3 shows how long it takes the
Sun to return to its maximum height at precisely noon everyday. This
shift from position 1 to position 3 is called a solar day.\label{Figure Sidereal}}
\end{figure}


A solar day is the amount of time it takes for the sun to return back
to the same position on Earth. Another important thing to keep track
of while trying to accurately track time is to count seconds. One
might think that counting days is more practical, but in fact, days
can be fairly complicated. Since the Earth is moving around the Sun
while it is spinning around it's axis, the Earth must actually rotate
a little more than 60\textdegree{}, or 2\textgreek{p} radians, in
order to return to the same position each day. Instead of spinning
precisely 360\textdegree{} every 24 hours, or 86400s, it actually
takes the Earth about 86164.0908s. This amount of time called a sidereal
day and is slightly longer than a solar day (see Figure \vref{Figure Sidereal}).
For every day of the year except January 1, sidereal and solar days
are out of sync. When trying to accurately track time, we really should
start from January 1. However, the question then becomes what time
and year relative to January 1 should one begin counting from?

In astronomy, astronomers typically count time in solar days starting
from noon, January 1, 4713 BC, also known as the Julian date. On the
other hand, computers count time in seconds from midnight at the start
of January 1, 1970. This is called \textquotedblleft{}Unix Time\textquotedblright{}
and is pretty close to Coordinated Universal Time (UTC) which is derived
directly from atomic clocks. Today, there have been a few revisions
to UTC with several leap seconds added in here and there. Additionally,
because the Earth's rotation is slowly slowing down, and keeping time
can become extremely complicated. By using Network Timing Protocol
(NTP) and/or the Global Positioning System (GPS), computers and cell
phones routinely synchronize their clocks to atomic standards with
milliseconds of precision. Because UTC is widely available to most
people today, this makes it relatively easy for us Astronomers to
accurately track time.


\subsection{Coordinate Conversion}


\subsubsection{Longitude and Latitude}

Longitude and latitude coordinates are considered geographic coordinates
(see Figure \vref{Latitude & Longitude Figure}). A geographic coordinate
system is a coordinate system that enables any location on the Earth
to be specified by a set of numbers or letters. The latitude of a
point on the Earth's surface is the angle between the equatorial plane
and the straight line that passes through that point normal to the
surface. Lines joining points of the same latitude form circles on
the surface of the Earth known as parallels because they are parallel
to the equator and to each other. While the +90\textdegree{} parallel
represents the north pole and -90\textdegree{} parallel represents
the south pole, the 0\textdegree{} parallel of latitude is at the
equator which divides the globe into Northern and Southern Hemispheres.

\begin{figure}[H]
\begin{centering}
\includegraphics[scale=0.48]{/Users/ppkantorski/Documents/Radio_Astronomy/Lab_3/Pictures/latitude_longitude}
\par\end{centering}

\caption{On Earth, latitude coordinates are measured as angles from the Equator
to +90� (North Pole) to -90� (South Pole). Longitude coordinates are
measured from the Prime Meridian from 0� to +180� due right and from
0 to -180� due left.\label{Latitude & Longitude Figure}}
\end{figure}


The longitude of a point on the Earth's surface is the angle east
or west from the reference meridian to another meridian that passes
through that point. The Prime Meridian is the line chosen as the international
zero-longitude reference line and was set with the intention of passing
through the British Royal Observatory in Greenwich, UK. Places to
the east of the Prime Meridian are considered to be in the eastern
hemisphere while places to the west are in the western hemisphere.
Longitude also ranges from about -180� to +180� relative to the Prime
Meridian.

Additionally, converting from rectangular latitude and longitude coordinates
to spherical coordinates can be done using the vector matrix shown
below in Equation \vref{eq:Long / Lat Rotation}.

\begin{equation}
\vec{x}=\left[\begin{array}{c}
\cos(lat)\cos(long)\\
\cos(lat)\sin(long)\\
\sin(lat)
\end{array}\right]\label{eq:Long / Lat Rotation}
\end{equation}



\subsubsection{Right Ascension and Declination ($\alpha,\sigma$)}

In spherical coordinates, the position of stars are often expressed
in right ascension and declination without use of a coordinate for
distance. Because celestial objects are extremely far way, astronomers
often have little to no information about their exact distances and
therefore use only the object's direction. Additionally, the direction
of fairly distant objects is the same for all observers and it is
convenient to specify this direction with the same coordinates for
every object. On the contrary, through the horizontal coordinate system,
a star's position differs from observer to observer based on their
positions on the Earth's surface and continuously changes with the
Earth's rotation.

\begin{figure}[H]
\begin{centering}
\includegraphics[scale=0.46]{/Users/ppkantorski/Documents/Radio_Astronomy/Lab_3/Pictures/ra_dec}
\par\end{centering}

\caption{The figure above shows right ascension and declination as seen on
the inside of the celestial sphere. The primary direction of the system
is the vernal equinox or where the ascending node of the ecliptic
in red meets the celestial equator in blue. Right ascension is measured
eastward along the celestial equator from the primary direction while
declination is measured due north or south from the equator.\label{RA & Dec Figure}}
\end{figure}


The equatorial coordinate system is a widely used celestial coordinate
system used to specify the positions of celestial objects. Two coordinates
typically used in the equatorial coordinate system are the right ascension
and declination (see Figure \vref{RA & Dec Figure}). Right ascension
(symbol \textgreek{a}, abbreviated RA) measures the angular distance
of an object eastward along the celestial equator from the vernal
equinox to the hour circle passing through the object. The vernal
equinox point is one of the two points where the ecliptic intersects
with the celestial equator. In comparison to longitude, right ascension
is similarly but is measured in sidereal hours, minutes and seconds
instead of degrees. This results from measuring right ascensions by
timing the passage of objects across the meridian as the Earth rotates.
In total, there are (360\textdegree{} / 24h) = 15\textdegree{} in
one hour of right ascension and 24h of right ascension around the
entire celestial equator.

The celestial equator is the projection of the Earth's equator onto
the celestial sphere. Declination (symbol \textgreek{d}, abbreviated
Dec) measures the angular distance of an object perpendicular to the
celestial equator. Declination ranges from +180� northward to -180�
to the south. The north celestial pole, for instance, has a declination
of +90\textdegree{}. The origin of the coordinate of declination is
the celestial equator itself. Declination can also be thought of as
analogous to latitude. When the two coordinates are used together,
right ascension and declination are typically abbreviated as (RA,
Dec) or ($\alpha,\sigma$). Because RA and Dec are an analogous to
longitude and latitude, Equation \vref{eq:Long / Lat Rotation} can
also be used to convert RA and Dec to spherical coordinates.


\subsubsection{Hour Angle and Declination ($ha,\sigma$)}

In the equatorial coordinate system, right ascension can also be translated
to hour angle (abbreviated HA or LHA). Hour angle positioning is a
``left-handed'' system which measures the angular distance of an
object westward along the celestial equator from the observer's meridian
to the hour circle passing through the object (See Figure \vref{Hour Angle Figure}).

\begin{figure}[H]
\begin{centering}
\includegraphics[scale=0.46]{/Users/ppkantorski/Documents/Radio_Astronomy/Lab_3/Pictures/hour_angle}
\par\end{centering}

\caption{The figure above shows right ascension measured eastward along the
celestial equator from the primary direction along with the hour angle
abbreviated by LHA. The hour angle coordinate is the right ascension
summed with the local sidereal time.\label{Hour Angle Figure}}
\end{figure}


Unlike right ascension, hour angle increases with the Earth's rotation.
Hour angle is used to measure the time from when an object crosses
the meridian. For instance, a star on the observer's celestial meridian
is said to have a zero hour angle. About one sidereal hour later,
the Earth's rotation will carry the star to the west of the meridian
and its hour angle becomes +1h. Right ascension may also be converted
into hour angle as an intermediate step to converting to azimuth and
altitude by multiplying the spherical coordinates with the rotational
matrix shown below in Equation \vref{eq:RA to HA}.

\begin{equation}
R_{(\alpha,\delta)\rightarrow(ha,\delta)}=\left[\begin{array}{ccc}
\cos(LST) & \sin(LST) & 0\\
-\sin(LST) & \cos(LST) & 0\\
0 & 0 & 1
\end{array}\right]\label{eq:RA to HA}
\end{equation}



\subsubsection{Azimuth and Altitude ($az,alt)$}

The horizontal coordinate system is a celestial coordinate system
that utilizes an observer's local horizon as it's fundamental plane.
This coordinate system divides the sky into an upper hemisphere and
a lower hemisphere. In the upper hemisphere, objects and stars are
visible or within sight while in the lower hemisphere, the our view
of the sky is blocked by the Earth. The ``great circle'' separating
hemispheres is called celestial horizon or rational horizon and the
pole of the upper hemisphere is called the zenith. 

Altitude (abbreviated Alt) is sometimes referred to as elevation or
the angle between the object and the observer's local horizon. For
visible objects in the sky, altitude is an angle between 0� to 90�.
Azimuth (abbreviated Az) is the angle of the object about the horizon
and is usually measured from the north while increasing towards the
east. Because the horizontal coordinate system is fixed to the Earth
and not the stars, the altitude and azimuth of an object changes with
time as the object appears to ``drift'' across the sky. Furthermore,
because the horizontal system is defined by the observer's position
or local horizon, different locations on Earth viewing the same object
at the same time will have different values of altitude and azimuth.

By multiplying spherical hour angle and declination coordinates by
the rotational matrix shown below in Equation \ref{eq:HA, Dec to Az, Alt},
we can convert ($ha,\sigma$) coordinates to ($az,alt)$ coordinates.
Since hour angle and declination are ``Earth-based'' coordinates,
this conversion will only depend on the observer's terrestrial latitude
\textgreek{f}. Additionally, we can also convert directly from right
ascension and declination to azimuth and altitude by multiplying our
coordinates by both rotational matrices.

\begin{equation}
R_{(ha,\delta)\rightarrow(az,alt)}=\left[\begin{array}{ccc}
-\sin(\phi) & 0 & \cos(\phi)\\
0 & -1 & 0\\
\cos(\phi) & 0 & \sin(\phi)
\end{array}\right]\label{eq:HA, Dec to Az, Alt}
\end{equation}



\subsection{Interferometry}

Interferometry is the practice of using two or more radio telescopes
in a synchronized array to observe astronomical sources and make data
measurements. A basic interferometer is composed of the array of radio
telescopes along with the electronics used to measure, record, and
detect signals. For this lab, my partners and I used a narrow band
two-element interferometer located on the roof of Kroeber Hall at
UC Berkeley, separated along an approximately East-West baseline.
Additionally, the interferometer operates at about 10.67 GHz, the
baseline separation between the two dishes is approximately B = 10-m,
and the observe radiation is at a wavelength of \textgreek{l} = 2.5
cm.


\subsubsection{The Two-Element Interferometer}

A two-element interferometer (see Figure \vref{two element inter Figure})
consists of two telescopes separated by a vector $\vec{b}$ known
as the baseline. In this system, both antennas simultaneously receive
electromagnetic radiation from a specified astronomical source located
above in the sky in the direction of the unit vector $\hat{s}$. Since
astronomical sources are extremely far away, radiation is transmitted
to the antennas in the form of ``plane waves.'' A plane wave is
a constant-frequency waveform in which the wavefronts are infinite
parallel planes of constant peak-to-peak amplitude normal to the phase
velocity vector.

\begin{equation}
\ensuremath{\vec{b}}\cdot\hat{s}=\left|\ensuremath{\vec{b}}\right|\cos\theta,\hspace{3em}\tau=\frac{\ensuremath{\vec{b}_{ij}}\cdot\hat{s}}{c}\label{eq:baseline}
\end{equation}


\begin{figure}[H]
\begin{centering}
\includegraphics[scale=0.75]{/Users/ppkantorski/Documents/Radio_Astronomy/Lab_3/Pictures/two_element_inter}
\par\end{centering}

\caption{When an incoming radio wave approaches the Earth at an angle, the
radio wave will arrive at the two separated telescopes at slightly
different times. In this picture, the delay in arrival time, also
known as the geometric delay, is the distance X divided by the velocity
of light. Also, the angular resolution of the interferometer is equal
to the wavelength divided by the effective baseline of the telescopes.
When the object being observed is directly overhead, the effective
baseline is approximately equal to the distance between the two telescopes.\label{two element inter Figure}}
\end{figure}


As the plane waves from astronomical sources reach the radio dishes,
one antenna will end up receiving the signal before the other. The
time that passes between the receiving of the signals is calls the
geometric delay and is denoted with the variable \textgreek{t}. In
order to accurately find the value of \textgreek{t}, we must know
the ``extra distance'' that the plane waves have to travel before
reaching the second telescope. This distance is approximately equal
to the baseline vector dotted with the directional unit vector. Additionally,
since the direction of the source is a unit vector, this distance
is approximately the projection of the baseline vector onto the unit
vector (see Equation \vref{eq:baseline}) with $\theta$ being the
angle between the baseline and direction vectors. If the baseline
projection of the interferometer is known, then the geometric delay
\textgreek{t} can be calculated by dividing this distance by the velocity
of the plane waves (the speed of light, c). When this delay is an
exact multiple of the wavelength, then the waves detected by the two
telescopes will be in phase and will add together when combined. If
they are not, then they will be out of phase and interfere with one
another. 

If there is another source in the sky at a different direction from
the first, the second source will also cause a geometric delay between
the two antennas. However, this would be different from the geometric
delay of the first source because it would have different directional
unit vectors $\hat{s}$. Because each antenna is receiving a signal
from each source, what is detected is actually the total signal of
the two. In order to separate the signals from each source at each
radio dish, we must correlate the total signals with each other.


\subsubsection{Correlation}

Correlation is very similar to convolution theorem. However, there
is a slight difference. When correlating two signals, the Fourier
transform of the second function is conjugated before being multiplied
and integrated over. Additionally, by correlating a flat-spectrum
signal with a time-shifted version of itself, this will yield a measure
of the power of the signal at the delay corresponding to the time
shift. For a two-element interferometer, the total signal detected
by each antenna can be represented as shown in Equation \vref{eq:total signal},
where $e{}_{1}$ and $e{}_{2}$ are the signals given by each source
at times t and the difference between t and either $\lyxmathsym{\textgreek{t}}_{1}$
or $\lyxmathsym{\textgreek{t}}_{2}$. 

\begin{equation}
E(t)=e_{1}(t)+e_{2}(t),\hspace{3em}E_{j}(t)=e_{1}(t-\tau_{1})+e_{2}(t-\tau_{2})\label{eq:total signal}
\end{equation}


\begin{equation}
(f\star g)(\tau)=\int{f(t)g^{*}(t-\tau)}dt\label{eq: correlation}
\end{equation}


By using $E_{i}(t)$ and $E_{j}(t)$ as $f$ and $g$ in the correlation
equation (see Equation \vref{eq: correlation}) respectively, the
correlation equation then becomes:

\begin{equation}
(f\star g)(\tau)=\int{e_{1}(t)e_{1}(t-\tau_{1}-\tau)+e_{2}(t)e_{2}(t-\tau_{2}-\tau)}dt\label{eq:correlation application 1}
\end{equation}


In order to extract meaningful signal information out of the integral
shown above in Equation \vref{eq:correlation application 1}, the
two functions must both not average to zero. This is only possible
when \textgreek{t} is equal either to $-\lyxmathsym{\textgreek{t}}_{1}$
or $-\lyxmathsym{\textgreek{t}}_{2}$. In the case of \textgreek{t}
= $-\lyxmathsym{\textgreek{t}}_{1}$, the correlation integral of
the two sources becomes:

\begin{equation}
(f\star g)(\tau)=\int e_{1}(t)e_{1}(t-\tau_{1}-(-\tau_{1}))dt=\int e_{1}(t)e_{1}(t)dt=\,<e_{1}^{2}>
\end{equation}


In this expression, the $e{}_{2}$ term averages to zero because $\tau_{2}-(-\tau_{1})$
does not cancel. The final result is the average power received by
the antennas for either the first or second source depending on the
geometric delay \textgreek{t}. So for a set of different geometric
delays, there is a set of power values that rises above signals that
cannot be correlated. The set of power values are essentially noise.


\subsubsection{Visibility Equation}

A single frequency $\nu$ with corresponding wavelength \textgreek{l}
has a number of wavelengths according to the relationship to the baseline
$\tau_{ij}\nu=\frac{\vec{b}\cdot\hat{s}}{\lambda}$. By knowing the
number of wavelengths between the two antennas, we can say that for
a signal of a particular frequency coming from a source in direction
$\hat{s}$, the complex phase difference between that signal measured
at antenna i and antenna j will be $e^{i\text{\textgreek{j}}}$, where
\textgreek{j} is the angle covered by the wave as it propagates from
i to j. Since \textgreek{j} is 2\textgreek{p} times the number of
wavelengths, we can the phase difference \textgreek{Df} (see Equation
\vref{eq:phase diff}).

\begin{equation}
\Delta\phi=e^{-2\pi i\tau_{ij}\nu}=e^{-2\pi i\frac{\vec{b}\cdot\hat{s}}{\lambda}}\label{eq:phase diff}
\end{equation}


The phase difference \textgreek{Df} is as shown in the equation above
frequency dependent. At particular frequencies, \textgreek{Df} also
varies with position across the sky. The pattern traced by the complex
phase difference on the sky is called the \textquotedblleft{}fringe
pattern\textquotedblright{} of the interferometer. To extrapolate
from a single baseline in a single direction a picture of how an entire
array might respond to the entire sky that falls within the primary
beam of the correlated telescopes, we must define coordinates representing
the length of a baseline in units of wavelength (see Equation \vref{eq:baseline in wavelength}).
Additionally, we must split the source direction vector $\hat{s}$
into its components.

\begin{equation}
\frac{\vec{b}}{\lambda}\equiv(u,v,w),\hspace{3em}\hat{s}\equiv(l,m,\sqrt{1-l^{2}-m^{2}})\label{eq:baseline in wavelength}
\end{equation}


In Equation \vref{eq:baseline in wavelength}, u represents the east-west
component of the baseline, v the north-south component, w the vertical
(up-down) component, l the east-west direction on the sky, and m the
north-south direction. The $\sqrt{1-l^{2}-m^{2}}$ component of $\hat{s}$
comes from the fact that we restrict our direction to be a unit vector.
Using these components, the response of a baseline, or \textquotedblleft{}visibility\textquotedblright{}
V, as a function of the u,v,w separation of the antennas, integrated
over all source intensity I on the sky as a function of l and m can
be written as follows:

\begin{equation}
V(u,v)=\int\int{A(l,m)\cdot I(l,m)\cdot e^{-2\pi i(ul+vm+w\sqrt{1-l^{2}-m^{2}})}dl\ dm}\label{eq:Visibility}
\end{equation}


Equation \vref{eq:Visibility} is the complete form of the \textquotedblleft{}visibility
equation\textquotedblright{} and is also sometimes referred to as
the \textquotedblleft{}measurement equation\textquotedblright{} of
an interferometer. Variable A(l,m) is the response of the primary
beam of the antennas as a function of it's direction on the sky. In
general, A(l,m) and I(l,m) are grouped together because the sky is
always perceived through the filter of the primary beam. The product
between the two is sometimes referred to as the \textquotedblleft{}perceived
intensity'' of the beam.

Equation \vref{eq:Visibility} can be easier understood if we make
a simplifying assumption known as the \textquotedblleft{}flat-sky\textquotedblright{}
approximation. The ``flat-sky'' approximation is either that w =
0 or that the primary beam A(l,m) is sufficiently small that $l,m\ll1$,
making the component$\sqrt{1-l^{2}-m^{2}}\approx1$. For either approximation,
the response of a baseline does not have to account for the fact that
the sky is a ``curved surface'' of a sphere. With this assumption,
the $e^{-2\pi iw\sqrt{1}}$ term is no longer a function of l and
m and can therefore be removed from the integral (see Equation \vref{eq:visibility fixed}).

\begin{equation}
V(u,v)=e^{-2\pi iw}\int\int{A(l,m)\cdot I(l,m)\cdot e^{-2\pi i(ul+vm)}dl\ dm}\label{eq:visibility fixed}
\end{equation}


This version of the visibility equation shown above states that when
the telescopes are phased to a \textquotedblleft{}phase center\textquotedblright{}
through a choice of a corresponding component $e^{\text{\textminus}2\text{\textgreek{p}}iw}$
with w being the baseline component along the direction toward the
phase center in units of wavelength, the visibility V(u,v) is becomes
just the Fourier Transform of the perceived sky (see Figure \vref{beam Figure}).

\begin{figure}[H]
\begin{centering}
\includegraphics[scale=0.56]{/Users/ppkantorski/Documents/Radio_Astronomy/Lab_3/Pictures/beam}
\par\end{centering}

\caption{The true image of the sky is shown on left and the true UV plane is
shown on the right.\label{beam Figure}}
\end{figure}


To better understand the fringe-pattern of a baseline on the sky,
we can equivalently think of taking an image of the sky in $l,m$
coordinates and then Fourier transform to create a \textquotedblleft{}uv-plane\textquotedblright{}
with coordinates being ``inverse angles.'' Because an inverse angle
is the same thing as a wavelength, the uv-plane will have coordinates
of u,v. Furthermore, for each pair of antennas, there will be two
samples; one at $u,v$ and another at $\lyxmathsym{\textminus}u,\lyxmathsym{\textminus}v$.

\begin{figure}[H]
\begin{centering}
\includegraphics[scale=0.56]{/Users/ppkantorski/Documents/Radio_Astronomy/Lab_3/Pictures/beam_dirty_1}
\par\end{centering}

\caption{The array sampling pattern in the uv-plane is shown on the right and
the associated synthesized beam pattern or \textquotedblleft{}dirty
beam\textquotedblright{} on the left.\label{Dirty Beam Figure-1}}
\end{figure}


Since the sky is real-valued, these two Fourier components are related
by a complex conjugate. This means that if we measure V(u,v) at $u,v$,
we will also have to measure $V^{*}(u,v)$ at $\lyxmathsym{\textminus}u,\lyxmathsym{\textminus}v$.
The sampling of the uv-plane is done by multiplying the true uv-plane
with the sampling pattern computed. This result is what we would get
if we took visibilities measured from an interferometer and then placed
each visibility V(u,v) at the corresponding $u,v$ (along with $\lyxmathsym{\textminus}u,\lyxmathsym{\textminus}v$)
coordinates of a matrix. Moreover, if you then take the inverse Fourier
Transform of this sampled uv-plane, you get a dirty image (see Figure
\vref{Dirty Beam Figure-2}).

\begin{figure}[H]
\begin{centering}
\includegraphics[scale=0.56]{/Users/ppkantorski/Documents/Radio_Astronomy/Lab_3/Pictures/beam_dirty_2}
\par\end{centering}

\caption{The array sampling pattern in the uv-plane is shown on the right while
the associated synthesized beam pattern or \textquotedblleft{}dirty
beam\textquotedblright{} is on the left.\label{Dirty Beam Figure-2}}
\end{figure}


Because the image shown in Figure \vref{Dirty Beam Figure-2} is fairly
degraded from the original, this is usually called a \textquotedblleft{}dirty
image\textquotedblright{}. Information is thus lost while sampling
the uv-plane. If we multiply the true uv-plane by our sampling function,
this is equivalent to convolving the true sky with the Fourier Transform
of our sampling function, also known as the ``dirty beam''. The
dirty beam is what is convolved with the true sky to yield our dirty
image. Additionally, it is possible to recover something resembling
the true sky by deconvolving the dirty image with the dirty beam (see
Figure \vref{Deconvolution figure}).

\begin{figure}[H]
\begin{centering}
\includegraphics[scale=0.56]{/Users/ppkantorski/Documents/Radio_Astronomy/Lab_3/Pictures/deconvolution}
\par\end{centering}

\caption{This figure shows the result of a \textquotedblleft{}cleaned\textquotedblright{}
dirty image. Note that the point sources are recovered but the extended
square in the middle is slightly worse as a result of using a ``point-source
model'' in the deconvolution process.\label{Deconvolution figure}}
\end{figure}


Deconvolution attempts to compensate for the information that is lost
by sampling part of the uv-plane by inserting prior information about
the sky. This information can be along the lines of just knowing that
the sky is a set of point-sources to knowing that we want the smoothest
sky that fits our data to the level of noise. Nevertheless, deconvolution
is not a well-posed issue until you decide precisely what prior information
you want to use for the sky.


\section{Methods and Observations}


\subsection{Using the Interferometer}

In this lab, the interferometer used in this lab has a relatively
short baseline and the fringe spacing is larger than the size of all
sources except for the Sun and Moon. Because of this, all of these
sources looked like \textquotedblleft{}point sources\textquotedblright{}
from which all the radiation appears to come from a single point like
a star in optical astronomy. The output of our interferometer is a
sinusoidal-like signal called the \textquotedblleft{}fringe\textquotedblright{}.
All information for the data collected using the interferometer is
stored in the frequency, amplitude, and phase of the fringe. Additionally,
if we initially know the baseline, the fringe properties become indicators
of the point-source's declination.

The interferometer we used in this lab is also called a multiplying
interferometer. A multiplying interferometer takes the signals from
the two telescopes and multiplies them together using a mixer. We
repetitively measured and recorded the time average over a suitable
interval. Using the cross product between the signals is important
because its average value is zero unless there if there is a source.
Thus, any detected signal has to come from the sky instead of from
the instrument itself. When we observe a source, the fringe amplitude
has no zero offset and the amplitude is therefore directly proportional
to the source's brightness or flux.


\subsection{Measuring Accurate Declinations}

Since our interferometer has a baseline of about 10-m, we can recover
a fringe spacing of about $\frac{\text{\textgreek{l}}}{B}\sim10\text{\textasciiacute}$.
From a horizon-to-horizon measurement, we were also able to measure
the declination approximately 100 times more accurately. Furthermore,
we can also achieve partial angular resolution of the Sun and Moon
along with their diameters to a fraction of a percent.

\begin{figure}[H]
\begin{centering}
\includegraphics[scale=0.54]{/Users/ppkantorski/Documents/Radio_Astronomy/Lab_3/Pictures/Angular_diameter}
\par\end{centering}

\caption{The figure above illustrates the concept of ``angular diameter''
in units of arc minutes, arc seconds, and degrees.\label{arc diam fig}}
\end{figure}



\subsection{Object Tracking in Python}

To track astronomical sources with the interferometer using the programming
language Python, my lab partners and I used the ``PyEphem'' and
``Radiolab'' Python libraries to position our telescope and monitor
time. By using UTC, we were able to figure out what sidereal time
it is in England at a longitude of 0� in hours, minutes, and seconds.
Since right ascension is measured in units of time, this told us what
part of the sky is overhead at the prime meridian. From our current
longitude at Berkeley (37.8732� N, -122.2573�), we figured out the
time offset we had add to get the current local sidereal time. The
LST tells us what part of the sky is directly overhead at our location.
By subtracting from the RA the LST of our astronomical sources, we
were able to get the hour angle as it appears from our location. Finally,
from using the latitude, HA, and the declination of our sources, we
were able to figure out precisely how and where to point the telescope.


\subsection{Point Sources}

While choosing point sources, my partners and I had to be careful.
Since the fringe frequency is dependent on $\cos\lyxmathsym{\textgreek{d}}$,
this means that we did not measure \textgreek{d} directly, but rather
$\cos\text{\ensuremath{\lyxmathsym{\textgreek{d}}}}$. Southern sources
presented a minor problem for horizon-to-horizon tracking because
of the presence of the Campanile at UC Berkeley. To compensate for
this, we had to delete some of our data. Additionally, our telescope
looking for Northern sources may be inaccurate, leading to loss of
the signal and noise. Our observation frequency was also in the TV
satellite band. Because geostationary TV satellites resides in the
southern skies, they can generate strong signals that could enter
the telescope's dish sidelines and produce fringes.

\begin{figure}[H]
\begin{centering}
\includegraphics[scale=0.58]{/Users/ppkantorski/Documents/Radio_Astronomy/Lab_3/Pictures/sources}
\par\end{centering}

\caption{Pictures of Orion Nebula, M17 Nebula, and Virgo A respectively.\label{sources Figure}}
\end{figure}


The three sources that my partners and I decided to observe were the
Orion Nebula, M17 Nebula, and Virgo A. From our observations, my partners
and I measured the approximate baseline of our interferometer using
fixed declination values for each source. This was done by using the
method of ``least-squares fitting'' in Python on the fringe phase
to the hour angle on fringe amplitude Equation \vref{eq:fringe amplitude for pt}
shown below. 

\begin{equation}
F\left(h_{s}\right)=A\cos\left[(2\pi(\frac{B_{y}}{\lambda}cos\left(\delta\right))sin\left(h_{s}\right)\right]-B\sin\left[(2\pi(\frac{B_{y}}{\lambda}cos\left(\delta\right))sin\left(h_{s}\right)\right]\label{eq:fringe amplitude for pt}
\end{equation}


\begin{figure}[H]
\begin{centering}
\includegraphics[scale=0.52]{/Users/ppkantorski/Documents/Radio_Astronomy/Lab_3/Pictures/untreated_orion}
\par\end{centering}

\caption{Untreated data from the Orion Nebula, collected from March 21, 2014
10:53PM - March 22, 2014 2:20AM.\label{Untreated Orion Figure}}
\end{figure}
\begin{figure}[H]
\begin{centering}
\includegraphics[scale=0.52]{/Users/ppkantorski/Documents/Radio_Astronomy/Lab_3/Pictures/untreated_M17}
\par\end{centering}

\caption{Untreated data from the M17 Nebula, collected from April 5, 2014 10:07AM
- 5:04PM.\label{Untreated M17 Figure}}
\end{figure}
\begin{figure}[H]
\begin{centering}
\includegraphics[scale=0.52]{/Users/ppkantorski/Documents/Radio_Astronomy/Lab_3/Pictures/untreated_virgo}
\par\end{centering}

\caption{Untreated data from Virgo A, collected from April 3, 2014 3:03AM -
4:20PM.\label{Untreated Virgo Figure}}
\end{figure}



\subsection{The Sun and Moon}

The Sun and Moon are special cases cases of astronomical sources for
two reasons. First of all, the position of the Sun and Moon change
from day to day. For the Moon, even a change of just an hour is extremely
significant. Secondly, the Sun and Moon are both extended sources,
not point sources. This causes their fringe amplitudes and phases
to change with time depending on their brightness distribution and
can make the determination of accurate positions somewhat complex.
Nevertheless, my partners and I were able to get the Sun and Moon
positions directly from the Python library ``PyEphem''.

\begin{figure}[H]
\begin{centering}
\includegraphics[scale=0.52]{/Users/ppkantorski/Documents/Radio_Astronomy/Lab_3/Pictures/untreated_sun}
\par\end{centering}

\caption{Untreated data from the Sun, collected from April 1, 2014 3:25PM -
April 2, 2014 1:02AM.\label{Untreated Sun Figure}}
\end{figure}
\begin{figure}[H]
\begin{centering}
\includegraphics[scale=0.46]{/Users/ppkantorski/Documents/Radio_Astronomy/Lab_3/Pictures/untreated_moon}
\par\end{centering}

\caption{Untreated data from the Moon, collected from April 4, 2014 6:54PM
- April 5, 2014 6:09AM.\label{Untreated Moon Figure}}
\end{figure}


Since the Earth's Moon it is relatively nearby, we also had to correct
for our location on the Earth\textquoteright{}s surface. The term
\char`\"{}parallax\char`\"{} refers to the semi-angle of inclination
between two sight-lines to the star as observed when the Earth is
on opposite sides of the Sun in its orbit. The parallax effect for
the Moon is far greater than that of the Sun. Unless if this is corrected
for, the Moon will probably lie outside the telescope's main beam.
Fortunately, PyEphem also automatically takes care of this along with
other complex issues like the ``time-of-flight'' of photons. Regardless
of these difficulties, the Sun and Moon are still good sources because
they are really bright and included them in our observations because
they produce huge signals which is ideal for testing our observations
and data analysis. Furthermore, the Sun is so bright that we got a
huge signal to noise ratio. This allowed us to estimate an accurate
declination.

From our observations, my partners and I measured the approximate
phase offset $\phi$ of our interferometer using the average$\frac{B_{y}}{\lambda}$
calculated from the point source data each source. This was done by
using the method of ``least-squares fitting'' in Python on the fringe
phase to the hour angle on fringe amplitude Equation \vref{eq:fringe amplitude for sun/moon}shown
below. 

\begin{equation}
R\left(h_{s}\right)=\cos\left(\frac{2\pi B_{y}}{\lambda}cos\left(\delta\right)\sin\left(h_{s}\right)+\phi\right)\label{eq:fringe amplitude for sun/moon}
\end{equation}



\subsection{Additional Details}

Since the two-element interferometer telescope has different distances
for each antenna from the source, it\textquoteright{}s convenient
to think of the different distances in terms of relative path delay
in units of time for the two telescopes. As stated earlier, this is
called the geometrical path delay$\tau_{g}$. However, the signals
have to travel through a lot of electronics before they are finally
multiplied. Additionally, these two paths are not necessarily of equal
length. This makes there an additional relative delay from the difference
in cable length $\tau_{c}$. The total relative delay is the sum of
the two (see Equation \vref{eq:relative delay}). While $\tau_{g}$
is a function of time, $\tau_{c}$ is independent of time.

\begin{equation}
\tau_{total}=\tau_{g}(h_{s})+\tau_{c},\hspace{3em}\tau_{g}(h_{s})=\left[\frac{B_{y}}{c}\cos\delta\right]\sin h_{s}\label{eq:relative delay}
\end{equation}


The local fringe frequency is the frequency of a certain point on
the sky. To see this, we can expand the hour angle term $\sin h_{s}$
into a Taylor series centered about the current hour angle of the
source $h_{s}$ (see Equation \vref{eq:expand hour}). The local fringe
frequency is contained in the second term of Equation \vref{eq:expand hour}
because for a small region around $h_{s}$ in Equation \vref{eq:fringe amplitude for pt},
$F(h)$ varies as $f_{f}\text{\ensuremath{\Delta}}h$. The local fringe
frequency $f_{f}$ is shown in Figure \vref{eq:fringe freq}.

\begin{equation}
\sin h=\left.\sin h_{s}+\triangle h\frac{d\sin h}{dh}\right|_{h_{s}}=\sin h_{s}+\triangle h\cos h_{s}\label{eq:expand hour}
\end{equation}


\begin{equation}
f_{f}=\frac{C}{2\pi}\cos(h_{s})=\left[\frac{B_{y}}{\lambda}\cos\delta\right]\cos h_{s}\label{eq:fringe freq}
\end{equation}


Equation \vref{eq:fringe freq} represents the local fringe frequency
in cycles per radian on the sky. To convert to cycles per hour, we
can multiply $\frac{dh_{s}}{dy}=\frac{2\pi}{24}$ where the number
24 represents the number of hours in a day. For instance, at the meridian,
assuming $B_{y}\approx10m$ and $\lyxmathsym{\textgreek{l}}=2.5cm$,
the fringe frequency is approximately $0.029\cos\text{\textgreek{d}}$
cycles per second while the period is $\frac{1}{f_{f}}=\frac{35}{\cos\delta}=35s$.


\subsubsection{Least Square Fitting}

A least squares method can be used on the filtered data from the point
sources and from the Sun and Moon using the fringe amplitude (Equation
\vref{eq:fringe amplitude for pt} and \vref{eq:fringe amplitude for sun/moon}).
For instance, for the point sources, this is done by minimizing the
square of the residuals ($\chi^{2}$) for particular values of $\tau_{c}$
and$B{}_{y}$. Using the following procedure shown below, my lab partners
and I were able to fit for both $\tau_{c}$ and$B{}_{y}$.

(1) Using python, we inserted Equation \vref{eq:fringe amplitude for pt}
with guessed constant coefficients for the$\tau_{c}$ dependent factors.

(2) Next, we iterate over a range of possible baseline values $B{}_{y}$,
and apply a least squares fit to each value. This provided us with
a $\chi^{2}$ and values of A and B.

(3) By looking at the specific value of $B_{y}$ corresponding to
the minimum $\chi^{2}$ value, this gives us the ``least squares
fit'' for $B_{y}$ , $A$, and $B$.


\subsubsection{Measuring The Diameter Of A Circular Source}

In this lab, goal was to measure and compare the diameters of the
Sun and Moon. Assuming that the sources are uniformly-bright disks
of radius R, this makes the intensity distribution approximately equal
to Equation \vref{eq:Intensity distribution circular}.

\begin{equation}
I(\triangle h)=\frac{(R^{2}-\triangle h^{2})^{1/2}}{R}\label{eq:Intensity distribution circular}
\end{equation}


Additionally, to obtain the theoretical modulating function $MF{}_{theory}$,
we can use the following expression:

\begin{equation}
MF_{theory}=\frac{1}{R}\intop_{-R}^{R}(R^{2}-\Delta h^{2})^{1/2}\cos(2\pi f_{f}\Delta h)d\Delta h\label{eq:modulation function}
\end{equation}


If we want to calculate $MF{}_{theory}$ analytically, this can be
done by substituting $\Delta(R\cos\theta)$ for $\Delta h$. The result
is a Bessel function. To do this, we can split $I(\Delta h)$ into
2N+1 tiny slices. Since the total number is odd, this makes the slices
symmetric about $\Delta h=0$. Each of these slices has width $\delta h=\frac{R}{N}$
along with $\Delta h_{n}=n\delta h$, where $n$ ranges from $-N$
to $+N$. After being rewritten, the integral becomes the sum shown
below.

\begin{equation}
MF_{theory}\approx\delta h\sum_{n=-N}^{n=+N}\left[1-(\frac{n}{N})^{2}\right]^{1/2}\cos(\frac{2\pi f_{f}Rn}{N})\label{eq:mod funct final}
\end{equation}



\section{Data Analysis}


\subsection{Python Coding to Clean Data}

To treat the data my partners and I gathered, I made a Python script
to eliminate bad data points for all my data, normalize the plot to
the origin, fit an envelope for the Sun and Moon data, plot the power
spectrum, and perform a ``Least Square Mean'' curve fit to return
the baseline $B_{y}$ for point sources, and use polynomial curve
fitting for the a parabolic portion of the Sun and Moon to derive
a Bessel function, and calculate the angular diameter of the sources.
Additionally, my script was also built to make plots for just about
any datasets taken from the interferometer given a specified file
path.

To normalize datasets to the x-axis, I using a method of data smoothing
called ``Boxcar Averaging''. This was done by averaging 100 points
above and below each datapoint to create an array of average data.
Initially, I used the Boxcar Average to find data points to cut out
of the raw data before entering the Boxcar average. Then, by subtracting
out the average data plot from the original graph, the data becomes
centered about roughly ``y = 0''. This ``cleaning'' my data data
calculations and analysis possible.

The next thing I did was fit a curve to the envelope of the Sun and
Moon's normalized data. To find the envelop, I cut out from the data
points points whose derivative was greater than a number relatively
close to 0. This allowed me to only plot points along the peaks in
the signal. Then, by taking the absolute value of the truncated data
and using another Boxcar average on the data points to smooth the
data, I was able to get a nice envelope for the data. For the Moon,
however, because the data points had a lot more spikes, I had to use
a scaling factor of about $2.7$ to scale my envelop to the correct
amplitude.

After finding the envelop curve for my Sun and Moon data, I was then
able to fit a parabolic polynomial curve to a ``parabola like''
region of my data. By using modulation function shown in Equation
\vref{eq:mod funct final} along with the constants taken from the
fitted polynomials, I was then able to derive a Bessel function to
approximate the diameter of the Sun and the Moon.


\subsection{Sun and Moon Data}


\subsubsection{The Sun}

\begin{figure}[H]
\begin{centering}
\includegraphics[scale=0.52]{/Users/ppkantorski/Documents/Radio_Astronomy/Lab_3/Pictures/treated_sun}
\par\end{centering}

\caption{Treated data with envelope and power spectrum from the Sun. Power
leakage occurred at about 25.71mHz.\label{Treated Sun Figure}}
\end{figure}
\begin{figure}[H]
\begin{centering}
\includegraphics[scale=0.52]{/Users/ppkantorski/Documents/Radio_Astronomy/Lab_3/Pictures/curve_sun}
\par\end{centering}

\caption{This figure shows a polynomial curve fit for the Sun data.\label{s_sq Sun Figure}}
\end{figure}
\begin{figure}[H]
\begin{centering}
\includegraphics[scale=0.52]{/Users/ppkantorski/Documents/Radio_Astronomy/Lab_3/Pictures/bessel_sun}
\par\end{centering}

\caption{This figure shows the Bessel function for our Sun data comparing $MF_{theory}$
from Equation \vref{eq:mod funct final} to $R\times f_{f}$ in order
to calculate the radius. Values used for $R\times f_{f}$ are when
the function above crosses the x-axis.\label{bessel sun fig}}
\end{figure}



\subsubsection{The Moon}

\begin{figure}[H]
\begin{centering}
\includegraphics[scale=0.52]{/Users/ppkantorski/Documents/Radio_Astronomy/Lab_3/Pictures/treated_moon}
\par\end{centering}

\caption{Treated data with envelope and power spectrum from the Moon. Power
leakage occurred at about 21.86.71mHz.\label{Treated Moon Figure}}
\end{figure}
\begin{figure}[H]
\begin{centering}
\includegraphics[scale=0.52]{/Users/ppkantorski/Documents/Radio_Astronomy/Lab_3/Pictures/curve_moon}
\par\end{centering}

\caption{This figure shows a polynomial curve fit for the Moon data.\label{poly moon fig}}
\end{figure}
\begin{figure}[H]
\begin{centering}
\includegraphics[scale=0.52]{/Users/ppkantorski/Documents/Radio_Astronomy/Lab_3/Pictures/bessel_moon}
\par\end{centering}

\caption{This figure shows the Bessel function for our Moon data comparing
$MF_{theory}$ to $R\times f_{f}$ in order to calculate the radius.
Values used for $R\times f_{f}$ are when the function above crosses
the x-axis.\label{bessel moon fig}}
\end{figure}



\subsubsection{Measuring the Diameter}

When measuring the diameter, because of my specific method of envelope
detection, it was not necessary to find $\phi$ from Equation \vref{eq:mod funct final}.
Instead, I ended up performing a nonlinear least square fitting to
fit a polynomial to my data and effectively return the Bessel function
associated to the data, the fringe frequency, and the approximated
radius. 

\begin{table}[H]
\begin{centering}
\begin{tabular}{|c|c|c|c|c|c|}
\hline 
Object Name: & HA & $\delta$ & $f_{f}$ (Hz) & Diameter & Actual Diameter\tabularnewline
\hline 
\hline 
Sun & -0.8865 & 0.7282 & 167.9125 & 30.9148' & 31.6\textasciiacute{} \textendash{} 32.7\textasciiacute{}\tabularnewline
\hline 
Moon & 1.4045 & 0.0000 & 58.9277 & 29.7526' & 29.3\textasciiacute{} \textendash{} 34.1\textasciiacute{}\tabularnewline
\hline 
\end{tabular}
\par\end{centering}

\caption{Hour angle, declination, fringe frequency, and source radius calculated
from Sun and Moon data.}
\label{diameter table}
\end{table}


As shown in the table above, the fringe frequencies and diameters
were close to what was to be expected. The Sun, however, was about
2.1684\% off from the average low. To approximate the right diameter
for the Sun, I used the second $R\times f_{f}$ value returned from
the Bessel function shown in Figure \vref{bessel sun fig}. For the
Moon, I used the first value shown on the Bessel function.


\subsection{Point Source Data}


\subsubsection{The Orion Nebula}

\begin{figure}[H]
\begin{centering}
\includegraphics[scale=0.56]{/Users/ppkantorski/Documents/Radio_Astronomy/Lab_3/Pictures/treated_orion}
\par\end{centering}

\caption{Treated data with power spectrum from the Orion Nebula. Power leakage
occurred at about 25.96mHz. \label{Untreated Orion Figure-1}}
\end{figure}
\begin{figure}[H]
\begin{centering}
\includegraphics[scale=0.52]{/Users/ppkantorski/Documents/Radio_Astronomy/Lab_3/Pictures/s_sq_orion}
\par\end{centering}

\caption{Calculation of $\frac{B_{y}}{\lambda}cos\left(\delta\right)$ for
the Orion using the ``Least Square Fitting'' method. Least squares
done on data ranged from about x = 2300s to 5000s. The dotted-line
shows what value made the best approximation to the baseline from
the data. \label{s_sq Orion Figure}}
\end{figure}



\subsubsection{M17 Nebula}

\begin{figure}[H]
\begin{centering}
\includegraphics[scale=0.56]{/Users/ppkantorski/Documents/Radio_Astronomy/Lab_3/Pictures/treated_M17}
\par\end{centering}

\caption{Treated data with power spectrum from the M17 Nebula. Power leakage
occurred at about 8.58mHz.\label{Untreated M17 Figure-1}}
\end{figure}
\begin{figure}[H]
\begin{centering}
\includegraphics[scale=0.52]{/Users/ppkantorski/Documents/Radio_Astronomy/Lab_3/Pictures/s_sq_M17}
\par\end{centering}

\caption{Calculation of $\frac{B_{y}}{\lambda}cos\left(\delta\right)$ for
the M17 using the ``Least Square Fitting'' method. Least squares
done on data ranged from about x = 3606s to 5604s. The two dotted-lines
acrossing shown in blue appeared to be fairly close to each other.
However, since the lowest $\chi^{2}$ value shown on the graph still
returned a closer approximation to the baseline, it was used for my
calculations. \label{s_sq M17 Figure}}
\end{figure}



\subsubsection{Virgo A}

\begin{figure}[H]
\begin{centering}
\includegraphics[scale=0.56]{/Users/ppkantorski/Documents/Radio_Astronomy/Lab_3/Pictures/treated_virgo}
\par\end{centering}

\caption{Treated data with power spectrum from Virgo A.\label{Untreated Virgo Figure-1}}
\end{figure}
\begin{figure}[H]
\begin{centering}
\includegraphics[scale=0.52]{/Users/ppkantorski/Documents/Radio_Astronomy/Lab_3/Pictures/s_sq_virgo}
\par\end{centering}

\caption{Calculation of $\frac{B_{y}}{\lambda}cos\left(\delta\right)$ for
Virgo A using the ``Least Square Fitting'' method. Least squares
done on data ranged from about x = 17123s to 22130s. Power leakage
occurred at about 7.28mHz.\label{s_sq Virgo Figure}}
\end{figure}



\subsubsection{Measuring the Baseline}

To measure an accurate value for the baseline, this can be done determined
by fitting interferometric data for a point source using Equation
\vref{eq:fringe amplitude for pt}. This equation has three unknown
quantities; the time delay due to the cable, $\tau_{c}$, the baseline
$B_{y}$, and the declination of the target $\lyxmathsym{\textgreek{d}}$.
To determine the baseline, my partners and I referred to the documented
declination of of our point sources seen in Table \vref{3C144 table}.
Table \vref{Baseline table} shown below has the derived coefficients
for Equation \vref{eq:fringe amplitude for pt} and baseline calculated
from using least square fitting on our point source data.

\begin{table}
\begin{centering}
\begin{tabular}{|c|c|c|c|}
\hline 
Object Name: & r.a.(2000) & dec(2000) & $S_{J_{y}}$\tabularnewline
\hline 
\hline 
Orion Nebula & $05^{h}35^{m}17.3s$ & $\lyxmathsym{\textminus}05^{\lyxmathsym{\textopenbullet}}23\lyxmathsym{\textasciiacute}28\lyxmathsym{\textasciiacute\textasciiacute}$ & $\sim340$\tabularnewline
\hline 
3C274 (Virgo A) & $12^{h}30^{m}49.423s$ & $+12^{\lyxmathsym{\textopenbullet}}23\lyxmathsym{\textasciiacute}28.04\lyxmathsym{\textasciiacute\textasciiacute}$ & $\sim34$\tabularnewline
\hline 
M17 & $18^{h}20^{m}26s$ & $\text{\textminus}16^{\text{\textopenbullet}}10.6\text{\textasciiacute}$ & $\sim500$\tabularnewline
\hline 
Sun & varies & varies & \tabularnewline
\hline 
Moon & varies & varies & \tabularnewline
\hline 
\end{tabular}
\par\end{centering}

\caption{3C144 Data Table for point sources and the Sun and Moon provided on
the UC Berkeley astronomy webpage AstroBaki.}
\label{3C144 table}
\end{table}
\begin{table}[H]
\begin{centering}
\begin{tabular}{|c|c|c|c|c|c|c|}
\hline 
Object Name & A & B & $\chi^{2}$ & Baseline$(B_{y})$ & $\Delta B_{y}$ & \% Error\tabularnewline
\hline 
\hline 
Orion & -2.1816e-06 & 2.4396e-06 & 2176.2744 & 9.7750 & 0.2250 & 2.2503 \tabularnewline
\hline 
M17 & 2.0591e-06 & -2.2544e-06 & 2164.1595 & 10.0765 & -0.0765 & 0.7651\tabularnewline
\hline 
Virgo A & 3.7002e-07  & 4.3418e-07 & 2258.3223 & 10.3395 & -0.3395 & 3.3948\tabularnewline
\hline 
\end{tabular}
\par\end{centering}

\caption{This table shows the resulting coefficients A and B from Equation
\vref{eq:fringe amplitude for pt} using least square fitting along
with the accuracy of fit or square of the residual $\chi^{2}$ . From
these results, the baseline information and relative error was derived.}
\label{Baseline table}
\end{table}


The average baseline was approximately 10.0637-m, a value within 0.6365\%
of the initial 10-m guess. Errors in data result may have resulted
from bad data measurements, possible (but not likely) flawes in coding,
external signal noise from the location of the telescopes, and/or
improper selection of data for the least square fitting.


\section{Conclusion}

Interferometry is an extremely useful tool for making more precise
measurements than from using single radio dishes. By synchronizing
two radio telescopes separated by a baseline, this effectively increases
the resolution of data measurements. Since there is a baseline separation
when using the interferometer, there is a distinct geometric delay
between the two dishes. This geometric delay is what produces the
``fringe patterns'' we observed. Fringe patterns can be used to
determine either the baseline of the interferometers or declination
of an astronomical point source being tracked.

When measuring the Sun or the Moon with the interferometer, the size
of the source is relatively larger than the interferometer we used.
Because of this, this gives rise to a modulating function. From the
modulating function, the angular radius can then be determined from
the function's shape. Although our data could have been better, results
were still fairly close to what we expected. Ultimately, this lab
has taught how to properly analyze singals and data in Python, use
an interferometer to make measurements, and aproximate the baseline
from point sources and angular diameter for the Sun and Moon. 


\section{Acknowledgements and References}
\begin{enumerate}
\item Lab Partners: Aaron Tran, Leonardo Cassara, and Ryan Gao.
\item Additional help with analysis, coding, and the report was also proved
by Baylee Bordwell and Eduardo Herrera.
\item Undergraduate radio lab. (2014, April 8). Retrieved from


https://casper.berkeley.edu/astrobaki/index.php/Undergraduate\_Radio\_Lab.

\item Additional pictures were taken from Wikipedia, the free online encyclopedia.\end{enumerate}

\end{document}
